\documentclass[letterpaper,11pt]{article}
\usepackage{amsmath, amssymb}
\usepackage{graphicx}
\usepackage{subcaption}
\usepackage{bridges}
\usepackage[colorlinks=true, urlcolor=blue, citecolor=black, linkcolor=black]{hyperref}
\setlength{\abovedisplayskip}{6pt}
\setlength{\belowdisplayskip}{6pt}

\urlstyle{rm}

% -----------------------------------------------------------------------------

\title{Resonator: Dynamic 3D Oscilloscope-Inspired Art}

\author{Guang Yang\\
University College Cork; gyang@umail.ucc.ie}

\date{}

% -----------------------------------------------------------------------------

\begin{document}

\maketitle
\thispagestyle{empty}

\begin{abstract}

\emph{Resonator} is an interactive generative artwork that extends oscilloscope imagery into three-dimensional space. Built in p5.js, the system uses harmonic motion, hypotrochoids, and torus knots to generate rotating networks of line segments. Viewer interaction modulates system behaviour, transforming deterministic equations into evolving spatial structures. By combining parametric geometry with real-time interaction, the work explores how simple oscillatory rules produce complex spatial forms.

\end{abstract}

%%%%%%%%%%%%%%%%%%%%%%%%%%%%%%%%%%%%%%%%%%%

\section*{Introduction}

Lissajous figures emerged in the nineteenth century through experiments by the French physicist Jules Antoine Lissajous, who used mirrors mounted on vibrating tuning forks to visualize harmonic relationships between oscillatory motions \cite{sierra2023recurrence}. Later, similar figures became closely associated with analog oscilloscopes operating in X--Y mode, where sinusoidal voltages applied to orthogonal axes generated dynamically evolving geometric traces. These images demonstrated that electronic signals could function simultaneously as analytical representations and as sources of visual abstraction.

Lissajous figures and related harmonic geometries have since appeared in scientific visualization, harmonographs, oscilloscope music, and mathematical art. In the 1950s, Ben Laposky used analog oscilloscopes to create his ``Oscillons,'' pioneering electronic abstractions generated from harmonic waveforms \cite{laposky1953electronic}. By feeding periodic motion into orthogonal coordinate systems, artists and researchers revealed intricate planar structures emerging from simple oscillatory rules \cite{grimsey2009lissajous,sierra2023recurrence}.

\emph{Resonator} is a web-based artwork implemented in p5.js, a JavaScript library for creative coding, that generates evolving three-dimensional structures in real time. The software samples a family of evenly spaced parameter values, evaluates paired parametric curves at each sample, and connects corresponding endpoint pairs with line segments. While contemporary visualizations of three-dimensional Lissajous figures often emphasize single continuous trajectories or isolated parametric curves \cite{wolfram3dlissajous2007}, \emph{Resonator} focuses on dynamically connected curve systems generated through paired parametric sampling and temporal modulation. The system contains three operating modes: Harmonic Phase, Hypotrochoidal Evolution, and Toroidal Knot Harmonics. The title refers to resonance metaphorically through recurring phase alignments and interference patterns rather than in the traditional physical sense. A live interactive version of the artwork is available online \cite{resonatorWeb}.

\begin{figure}[t]
\centering

\begin{subfigure}{0.32\linewidth}
\centering
\includegraphics[width=\linewidth]{resonator_1.png}
\caption{Harmonic Phase}
\end{subfigure}
\hfill
\begin{subfigure}{0.32\linewidth}
\centering
\includegraphics[width=\linewidth]{resonator_2.png}
\caption{Hypotrochoidal Evolution}
\end{subfigure}
\hfill
\begin{subfigure}{0.32\linewidth}
\centering
\includegraphics[width=\linewidth]{resonator_4.png}
\caption{Toroidal Knot Harmonics}
\end{subfigure}

\caption{Three generative modes of \emph{Resonator}.}
\end{figure}

%%%%%%%%%%%%%%%%%%%%%%%%%%%%%%%%%%%%%%%%%%%
\section*{Mathematical Structure}

\subsection*{Harmonic Phase}

The Harmonic Phase mode is built from layered harmonic motion. For each sample value $u_i=i/N$, the system computes two time-dependent phase angles,
\[
\theta_i(t)=2\pi f u_i+t, \qquad
\phi_i(t)=2\pi u_i+\frac{t}{2},
\]
and uses them to generate a pair of three-dimensional endpoints:
\[
\begin{aligned}
x_1 &= R \sin(\theta_i)\cos(m_\phi \phi_i), &
y_1 &= R \cos(m_\theta \theta_i)\sin(\phi_i), &
z_1 &= R \sin(\phi_i + \theta_i), \\
x_2 &= R \cos(\theta_i + t), &
y_2 &= R \sin(\phi_i + \theta_i), &
z_2 &= R \cos(\phi_i - \theta_i).
\end{aligned}
\]
Here $R$ is a scale radius, $f$ is the frequency multiplier, and $m_\phi$ and $m_\theta$ are additional harmonic multipliers. The explicit dependence on $t$ ensures that the endpoints evolve continuously over time.

This structure generalizes the classical Lissajous figure
\[
 x = \sin(at), \quad y = \sin(bt + \delta),
\]
where $a$ and $b$ are frequency multipliers and $\delta$ is a phase offset. Instead of tracing a single planar curve, the system distributes many phase-shifted instances across the parameter interval $u \in [0,1]$. For each sample value $u_i$, two spatial points are generated and connected by a line segment. The accumulation of these segments forms a rotating bundle of oscillatory vectors with recurring phase alignments.

All three modes share a common construction principle illustrated in Figure~2. Two parametric curves are sampled at shared parameter values $u_i=i/N$, and corresponding points are connected by line segments. Figure~2 shows planar projections of these pairings prior to three-dimensional rendering. Open circles indicate sample points on one curve, filled circles indicate the corresponding points on the second curve, and the green segment highlights one representative pairing.

\begin{figure}[h]
\centering
\begin{subfigure}{0.32\linewidth}
\centering
\includegraphics[width=\linewidth]{figure2a.png}
\caption{Harmonic Phase}
\end{subfigure}
\hfill
\begin{subfigure}{0.32\linewidth}
\centering
\includegraphics[width=\linewidth]{figure2b.png}
\caption{Hypotrochoidal Evolution}
\end{subfigure}
\hfill
\begin{subfigure}{0.32\linewidth}
\centering
\includegraphics[width=\linewidth]{figure2c.png}
\caption{Toroidal Knot Harmonics}
\end{subfigure}

\caption{Construction diagrams for the three generative modes of \emph{Resonator}. Curve A (open circles) and Curve B (filled circles) are sampled at shared parameter values and connected by line segments. The green segment highlights one representative pairing of corresponding points.}

\end{figure}

%%%%%%%%%%%%%%%%%%%%%%%%%%%%%%%%%%%%%%%%%%%
\subsection*{Hypotrochoidal Evolution}

In the second mode, the endpoints follow hypotrochoidal equations. A hypotrochoid is the curve traced by a point on a circle rolling inside a larger fixed circle. The standard parameters $R$, $r$, and $d$ denote the fixed-circle radius, the rolling-circle radius, and the pen offset:
\[
x = (R-r)\cos t + d\cos\left(\frac{R-r}{r}t\right), \qquad
y = (R-r)\sin t - d\sin\left(\frac{R-r}{r}t\right).
\]
For each sample value $u_i=i/N$, let
\[
\theta_i = 2\pi \kappa u_i, \qquad
t_{1,i}(t)=\theta_i+t, \qquad
t_{2,i}(t)=\theta_i-\lambda t,
\]
where $\kappa$ is the angular repetition factor and $\lambda$ controls the relative temporal drift between the two curves. The three-dimensional version adds sinusoidal $z$-components so that the planar roulette becomes a three-dimensional structure:
\[
z_{1,i}(t)=A_1\sin(a_1\theta_i+b_1 t), \qquad
z_{2,i}(t)=A_2\cos(a_2\theta_i-b_2 t).
\]
Two hypotrochoids with distinct parameter sets are paired at each sample value $u_i$, so every rendered segment joins corresponding points on related curves. The two curves use slightly different time offsets, which causes the line segments to continuously transform through space. Additive color blending enhances interference patterns and layered transparencies.

%%%%%%%%%%%%%%%%%%%%%%%%%%%%%%%%%%%%%%%%%%%

\subsection*{Toroidal Knot Harmonics}

The third mode uses torus-knot parameterizations and emphasizes interacting closed trajectories within a rotating frame \cite{sequin2012topological}. Here the sample parameter is written as $s_i=i/N$, reserving $\theta$ for the harmonic mode, with $\psi_i = 2\pi \kappa s_i$. Distinct integer pairs $(m,n)$ are used for each endpoint:
\[
u_{1,i}(t)=m_1\psi_i+t, \qquad
v_{1,i}(t)=n_1\psi_i-\alpha t,
\]
\[
u_{2,i}(t)=m_2\psi_i-\beta t, \qquad
v_{2,i}(t)=n_2\psi_i+\gamma t.
\]
The corresponding endpoints are
\[
x_1 = (R + r\cos(v_{1,i}))\cos(u_{1,i}), \qquad
y_1 = (R + r\cos(v_{1,i}))\sin(u_{1,i}), \qquad
z_1 = r\sin(v_{1,i}),
\]
\[
x_2 = (R_2 + r_2\cos(v_{2,i}))\cos(u_{2,i}), \qquad
y_2 = (R_2 + r_2\cos(v_{2,i}))\sin(u_{2,i}), \qquad
z_2 = r_2\sin(v_{2,i}).
\]
In the implementation, pairs such as $(3,2)$ and $(5,4)$ generate related but distinct trajectories.

Because $m$ and $n$ are integers, the resulting paths close periodically and wrap the surface of a mathematical torus \cite{oberti2016torus}. Their superposition generates transient braided or woven three-dimensional configurations before the periodic structure becomes perceptible.

%%%%%%%%%%%%%%%%%%%%%%%%%%%%%%%%%%%%%%%%%%%
\section*{Temporal and Chromatic Modulation}

Time enters the system as a continuously increasing phase parameter. Cursor distance from the center of the screen modulates amplitude and temporal speed. Radial distance scales both the base radius and phase advancement. Small gestures therefore expand or contract the entire spatial configuration.

Color is computed in HSB (Hue, Saturation, Brightness) space. A primary hue is selected and combined with a secondary hue offset. The final color of each segment is calculated from a sinusoidal function of time and segment index:
\[
H = \sin(\omega_t t + \omega_i i),
\]
where $i$ denotes the position of a segment within the parameterized sequence, $t$ is time, $\omega_t$ controls hue variation over time, and $\omega_i$ controls hue variation across the index of the segment. The resulting value is mapped between the two hue values. Because hue evolves simultaneously across time and index, waves of color propagate across the structure. Additive blending causes overlapping segments to accumulate brightness, reinforcing regions of geometric density. Motion and color therefore form parallel layers of periodic variation, reinforcing the sense that the structure behaves like a resonant visual system.

%%%%%%%%%%%%%%%%%%%%%%%%%%%%%%%%%%%%%%%%%%%
\section*{Interactivity and Generative Systems}

A defining feature of \emph{Resonator} is interactivity. Viewer input alters amplitude and temporal speed, while a click reverses the direction of global rotation. These actions do not interrupt the underlying equations; instead they adjust parameters within the generative system.

The software continuously rotates the structure about multiple axes, producing slow spatial drift even in the absence of interaction. Because each rendering mode contains hundreds of independently computed segments, small changes propagate across the entire structure.

This approach reflects principles of generative systems. A small set of rules based on harmonic oscillation and parametric rotation produces complex visual behavior. Although the equations remain fixed, the perceptual result evolves continuously through interaction.

Unlike historical oscilloscope images, which recorded fixed electronic signals, \emph{Resonator} remains in flux. It does not display a single curve but sustains a dynamic field of relationships. The viewer becomes part of the system’s feedback loop, influencing resonance without fully determining it.

%%%%%%%%%%%%%%%%%%%%%%%%%%%%%%%%%%%%%%%%%%%
\section*{Visual and Conceptual Implications}

The work reframes Lissajous figures, hypotrochoids, and torus knots as evolving structures rather than static diagrams. Through motion, overlap, and additive color blending, analytic curves become dynamic visual systems. In this sense, \emph{Resonator} aligns with forms of mathematical art in which parametric and topological forms operate simultaneously as mathematical models and aesthetic objects \cite{sequin2012topological}.

In summary, \emph{Resonator} extends oscilloscope-inspired imagery into an interactive three-dimensional generative system. The work transforms oscillatory equations into evolving three-dimensional structures. By combining deterministic geometry with real-time interaction, the project demonstrates how mathematical systems can produce rich visual complexity.

%%%%%%%%%%%%%%%%%%%%%%%%%%%%%%%%%%%%%%%
{\setlength{\baselineskip}{12pt}
\raggedright
\bibliographystyle{bridges}
\bibliography{references}
}

\end{document}
